\section{Class-Utility Augmentation}
\label{sec:method}

CUA decides, for every training image, which single edit to apply on top of a fixed base recipe. It makes this decision by keeping an estimate of how much each edit currently helps each class learn, and by combining these estimates over the classes present in the image. Fig.~\ref{fig:framework} shows the two parts of the method: an estimator that periodically updates a table of class utilities from gradients (Sec.~\ref{sec:method_estimator}), and an image policy that reads this table to choose the edit for each image (Sec.~\ref{sec:method_policy}). We first fix the notation and the augmentation that every image receives regardless of the policy.

\begin{figure*}[t]
\centering
\begin{tikzpicture}[
  font=\footnotesize, >={Stealth[length=5pt]}, node distance=5mm and 7mm,
  box/.style={draw, rounded corners=2pt, align=center, inner sep=3pt, minimum height=11mm},
  data/.style={box, fill=gray!10},
  key/.style={box, fill=gray!25, line width=0.9pt}]
  % training loop
  \node[data, text width=2.2cm] (img) {Training image $x$\\ classes $S(x)$};
  \node[box, right=of img, text width=2.2cm] (base) {Base recipe\\ (flip, scale jitter)};
  \node[key, right=of base, text width=4.0cm] (pol) {\textbf{Image policy}\\ score $s(o\mid S)$ over the classes in $x$\\ harm block, softmax sampling\\ ``no edit'' always allowed};
  \node[box, right=of pol, text width=2.0cm] (edit) {Apply edit $T_o$\\ $o\in\mathcal{O}\cup\{o_0\}$};
  \node[box, right=of edit, text width=2.4cm] (upd) {Detector update\\ $\theta\leftarrow\theta-\eta\,\nabla_\theta\mathcal{L}$};
  \draw[->] (img) -- (base);
  \draw[->] (base) -- (pol);
  \draw[->] (pol) -- node[above]{$o$} (edit);
  \draw[->] (edit) -- (upd);
  % estimator
  \node[data, below=19mm of img, text width=2.3cm] (B) {Estimation batch $B$\\ (8 train.\ images)};
  \node[box, right=6mm of B, text width=2.6cm, font=\scriptsize] (g) {$g_0=\nabla_{\theta_s}\mathcal{L}(B)$\\[2pt] $g_o=\nabla_{\theta_s}\mathcal{L}(T_o(B))$};
  \node[data, below=5mm of B, text width=2.3cm] (R) {Reference set\\ subsets $R_c$};
  \node[box, right=6mm of R, text width=2.6cm, font=\scriptsize] (h) {$h_c=\nabla_{\theta_s}\mathcal{L}_c(R_c)$};
  \coordinate (mid) at ($(g)!0.5!(h)$);
  \coordinate (eu) at ($(edit.center)!0.5!(upd.center)$);
  \node[box, text width=3.5cm] (u) at (pol |- mid) {\textbf{Class utility}\\[2pt] $u(o,c)=\dfrac{\eta\,\langle g_o-g_0,\,h_c\rangle}{\mathcal{L}_c(R_c)+\delta}$};
  \node[key, text width=4.4cm] (tab) at (eu |- mid) {\textbf{Utility table} $\bar u(o,c)$, $\mathrm{se}(o,c)$\\ edits $\times$ classes, moving average};
  \node[box, dashed, below=4mm of tab, text width=4.4cm] (sh) {Shared-utility variant: one pooled row\\ $\bar u_{\mathrm{sh}}(o)=\frac{1}{|\mathcal{C}|}\sum_{c}\bar u(o,c)$ for every image};
  \draw[->] (g.east) -- (u.west |- g.east);
  \draw[->] (h.east) -- (u.west |- h.east);
  \draw[->] (u) -- (tab);
  \draw[->, dashed] (tab) -- (sh);
  \draw[->, line width=0.9pt] (tab.north) -- ++(0,5mm) -| node[pos=0.25, above]{utilities} (pol.south);
  \draw[->, dashed] (img.south) -- node[right, font=\scriptsize]{sampled} (B.north);
  \begin{scope}[on background layer]
    \node[draw, dashed, rounded corners=3pt, fit=(B)(R)(g)(h)(u)(tab)(sh), inner sep=5pt] (frame) {};
  \end{scope}
  \node[anchor=south east, font=\footnotesize\itshape, inner sep=2pt] at (frame.north east) {class-utility estimator, every $K$ optimiser steps};
\end{tikzpicture}
\caption{Overview of Class-Utility Augmentation (CUA). Top: in every training step, each image receives the base recipe and then one edit, chosen by the image policy from the classes the image contains; ``no edit'' is always a candidate. Bottom: every $K$ steps, the estimator compares the gradients of the training loss on an edited and an unedited estimation batch with the gradient of each class's loss on its reference subset. Their alignment estimates how much each edit helps each class learn (Eq.~\eqref{eq:utility}), and the smoothed estimates form the utility table that the policy reads. Shared-utility augmentation (dashed) pools the table into one row, so every image receives the same preferences; it differs from CUA in nothing else.}
\label{fig:framework}
\end{figure*}

\subsection{Problem Setup}
\label{sec:method_setup}

Let $f_\theta$ be a detector with parameters $\theta$, trained on images annotated with boxes from a set of classes $\mathcal{C}$. A training image $x$ carries annotations $Y=\{(b_j,c_j)\}_{j=1}^{n}$, and we write $S(x)=\{c_1,\dots,c_n\}\subseteq\mathcal{C}$ for the set of classes it contains, so that $|S(x)|\geq 2$ for a multi-class image. Training minimizes the standard detection loss $\mathcal{L}(\cdot\,;\theta)$, which sums classification and box-regression terms over the matched proposals or anchors.

Every training image first receives a base recipe that all compared policies share: a horizontal flip with probability 0.5, followed by isotropic scale jitter with a factor drawn from $[0.8,1.2]$ on a $640\times640$ letterboxed canvas, where a draw that pushes an annotated box out of the canvas is redrawn. The base recipe is never adapted. On top of it, an augmentation policy chooses one edit $o$ from the candidate set $\mathcal{O}\cup\{o_0\}$, where $\mathcal{O}$ is the augmentation menu of the dataset (Table~\ref{tab:menu}) and $o_0$ is the identity, or ``no edit''. We write $T_o$ for the transformation that applies edit $o$ to an image together with its boxes. The menu covers geometric, colour, image-quality, occlusion and composition edits, each at one moderate strength, so the policy controls which edit an image receives rather than how strong that edit is. The task of the policy is to choose $o$ for each image so that the resulting training step benefits the classes the image contains.

\subsection{Estimating Class Utility}
\label{sec:method_estimator}

The quantity CUA needs is the effect that training on an edited image has on what the detector learns about each class. Measuring this effect directly requires training a separate detector for each edit, which is the costly protocol of Section~\ref{sec:measurement}; CUA instead estimates it during training from a single gradient step. Consider one stochastic gradient step with learning rate $\eta$ on a batch $B$ of training images. If the batch is edited by $o$, the step changes the parameters by $-\eta\,g_o$ with $g_o=\nabla_\theta\mathcal{L}(T_o(B);\theta)$, and if it is left unedited, by $-\eta\,g_0$ with $g_0=\nabla_\theta\mathcal{L}(B;\theta)$. Let $\mathcal{L}_c(R_c;\theta)$ be the loss of class $c$ on a held-out set $R_c$ of images that contain $c$, restricted to the proposals or anchors matched to class-$c$ ground truth, and let $h_c=\nabla_\theta\mathcal{L}_c(R_c;\theta)$. To first order, the step on the edited batch reduces this loss by $\eta\langle g_o,h_c\rangle$ and the step on the unedited batch by $\eta\langle g_0,h_c\rangle$. We define the class utility of edit $o$ for class $c$ as the difference between the two reductions, normalized by the current loss of the class:
\begin{equation}
u(o,c)=\frac{\eta\,\langle g_o-g_0,\;h_c\rangle}{\mathcal{L}_c(R_c;\theta)+\delta},
\label{eq:utility}
\end{equation}
with $\delta=0.01$. A positive $u(o,c)$ means that a step on the edited batch moves the detector further toward a lower loss for class $c$ than a step on the unedited batch would, so the edit currently teaches the class something useful; a negative value means that the edit pulls the learning of $c$ in the wrong direction. By construction, $u(o_0,c)=0$ for every class. The normalization by $\mathcal{L}_c$ makes utilities comparable across classes whose losses differ in scale, so that no class dominates the score of an image merely because its loss is large. Eq.~\eqref{eq:utility} is the first-order influence estimate of TracIn~\cite{pruthi2020tracin}, applied here to a pair of edited and unedited batches rather than to one training example, and to a separate reference loss for each class rather than to one test loss. Like TracIn, it ignores optimizer momentum and higher-order terms, and it predicts the effect of one step rather than of training to convergence; Section~\ref{sec:measurement} provides the measurements against which this approximation is checked.

The held-out images come from a \emph{reference set} of about 1,000 training images that are removed from optimization in every compared policy and are used only by the estimator, so that the utilities reflect images the detector has not trained on. For each class, $R_c$ is a subsample of 64 reference images that contain $c$. Every $K$ optimizer steps, CUA draws an estimation batch $B$ of eight training images, computes $g_0$ and $g_o$ for every valid edit using the same base-recipe draw for all of them, computes $h_c$ for every class, and evaluates Eq.~\eqref{eq:utility} for all pairs $(o,c)$. To keep this step affordable, all gradients are taken with respect to a parameter subset $\theta_s$ rather than the full network: the box head and the top feature-pyramid level for Faster R-CNN, and the classification and regression subnets for \todo{second detector}. These parameters are closest to the class decision, which is where a class-specific effect of an edit is expressed.

A single estimate is noisy, because it rests on eight training images and one draw of each edit. CUA therefore smooths the estimates with an exponential moving average,
\begin{equation}
\bar u(o,c)\leftarrow\beta\,\bar u(o,c)+(1-\beta)\,u(o,c),\qquad\beta=0.9,
\label{eq:ema}
\end{equation}
and tracks the uncertainty of the average with a standard error $\mathrm{se}(o,c)$ derived from an exponentially weighted variance of the instantaneous estimates \todo{exact se formula; confirm against the implementation}. Estimation starts after the first training epoch; until then no utilities exist, and the policy samples uniformly from $\mathcal{O}\cup\{o_0\}$. Because the table is refreshed throughout training, the edit that a class prefers can change as the detector improves.

\subsection{Choosing an Edit for Each Image}
\label{sec:method_policy}

The image policy turns the utility table into one decision per image. For an image with class set $S=S(x)$, CUA scores each candidate edit by the average utility of the classes present,
\begin{equation}
s(o\mid S)=\frac{1}{|S|}\sum_{c\in S}w_c\,\bar u(o,c),
\label{eq:score}
\end{equation}
with equal class weights $w_c=1$. Averaging over the classes in the image, rather than over all classes, is what makes the decision depend on the image: an image that contains only potholes is scored by the pothole entries of the table, whereas an image that contains a crack and a pothole is scored by both. Since the identity scores $s(o_0\mid S)=0$, an edit is preferred to leaving the image unchanged only when its average utility over the classes present is positive.

An average can still hide a loss for one of the classes, which is precisely the failure CUA is designed to avoid. Before sampling, CUA therefore removes every edit that is reliably harmful to at least one class in the image, and keeps the allowed set
\begin{equation}
\mathcal{A}(S)=\Bigl\{o:\ \min_{c\in S}\bigl[\bar u(o,c)+1.64\,\mathrm{se}(o,c)\bigr]\geq 0\Bigr\},
\label{eq:harm}
\end{equation}
where $o$ ranges over $\mathcal{O}\cup\{o_0\}$. The set excludes an edit whenever its utility for a present class is negative at one-sided 95\% confidence. The identity is never excluded, because its utility and standard error are zero, so $\mathcal{A}(S)$ is never empty. CUA then samples the edit from
\begin{equation}
p(o\mid S)=(1-\varepsilon)\,\frac{\exp\bigl(s(o\mid S)/\tau\bigr)}{Z(S)}+\frac{\varepsilon}{|\mathcal{A}(S)|},
\label{eq:sampling}
\end{equation}
for $o\in\mathcal{A}(S)$, where $Z(S)=\sum_{o'\in\mathcal{A}(S)}\exp\bigl(s(o'\mid S)/\tau\bigr)$ normalizes the softmax, $\varepsilon=0.1$, and the temperature $\tau$ is the standard deviation of the current scores over $\mathcal{A}(S)$, floored at $10^{-6}$. Tying $\tau$ to the spread of the scores keeps the distribution unchanged when all utilities shrink or grow together, as they do when the loss falls during training. The uniform term keeps some diversity among the allowed edits, which limits the damage when an estimate is wrong.

Two cases complete the policy. Images that contain no annotated object have an empty class set; in every compared policy, such an image receives either the identity or a uniformly sampled edit, each with probability 0.5, so the treatment of background-only images does not differ between methods. Mosaic combines four images; CUA chooses it from the class set of the anchor image, and the other three tiles are drawn uniformly from the training set.

\subsection{Shared-Utility Augmentation}
\label{sec:method_shared}

To isolate the contribution of class information, we define a counterpart that differs from CUA in that respect alone. Shared-utility augmentation runs the same estimator and the same policy, but it pools the instantaneous estimates over all classes, $u_{\mathrm{sh}}(o)=\frac{1}{|\mathcal{C}|}\sum_{c\in\mathcal{C}}u(o,c)$, smooths them as in Eq.~\eqref{eq:ema}, and applies the harm block of Eq.~\eqref{eq:harm} to the pooled estimate and its standard error. Every image therefore receives the same edit preferences, whichever classes it contains. Because the estimator, reference set, sampling rule and settings are shared, a difference in accuracy between CUA and shared-utility augmentation can be attributed to the use of per-class utilities.

\subsection{Computational Cost}
\label{sec:method_cost}

CUA adds no policy network, no search phase before training and no training images; its only cost is the periodic estimation step. Each estimation step evaluates $|\mathcal{O}|+1$ gradients on the eight-image estimation batch and $|\mathcal{C}|$ gradients on 64-image reference subsamples, all restricted to $\theta_s$, and its cost is amortized over the $K$ optimizer steps it serves. The cost therefore grows linearly with the size of the menu and with the number of classes. On RDD2022, with 12 edits and four classes, one estimation step processes 104 training and 256 reference images, against the 800 training images of the $K=200$ steps it serves; on PASCAL VOC, with 20 classes, the reference term grows to 1,280 images. We report the measured wall-clock overhead for every dataset together with the results.

\section{Measuring Class-Specific Augmentation Effects}
\label{sec:measurement}

CUA can help only if classes respond differently to the same edit, and its estimates are useful only if they track how classes actually respond. We therefore measure these responses directly, with controlled training runs that do not involve CUA. The measurements serve three purposes: they test whether augmentation effects in detection depend on the class, they provide the reference against which the class utilities of Section~\ref{sec:method} are checked, and they define the measured-effect policy used as a baseline in Section~\ref{sec:setup}. All measurements use the datasets, detectors and training recipe of Section~\ref{sec:setup}, and all are evaluated on the validation split.

\subsection{Edit-Wise Measurement Study}
\label{sec:measure_editwise}

The edit-wise study measures the effect of each edit on each class by training with one edit at a time. Training a full detector for every edit, dataset, detector and seed would be prohibitively expensive, so the study branches short training runs from a shared starting point. For each dataset, detector and seed, we first train a warm checkpoint with the base recipe for three epochs. From this checkpoint, we continue training for three further epochs in $|\mathcal{O}|+1$ branches: one with the base recipe alone, and one for each valid edit $o$, in which $o$ is applied to the whole image with probability 0.5 after the base recipe. The learning rate remains at its base value for the first two continuation epochs and is reduced tenfold for the last. All branches share the warm checkpoint and its optimizer state, the order of the training images, the base-recipe draws and the random draw that decides whether an edit is applied, so any two branches differ only in the edit itself. For edit $o$, class $c$ and seed $k$, the measured effect is the change in validation AP50 at the end of the continuation,
\begin{equation}
\Delta_{o,c,k}=\mathrm{AP50}_c(o,k)-\mathrm{AP50}_c(\mathrm{base},k),
\label{eq:delta}
\end{equation}
and averaging over seeds gives the edit--class effect table $\bar\Delta(o,c)$. The study covers the three datasets, both detector families and seeds $\{0,1,2\}$, for a total of 258 training runs including the warm checkpoints.

\subsection{Testing Whether Effects Depend on the Class}
\label{sec:measure_tests}

A difference between two classes in Eq.~\eqref{eq:delta} can arise from training noise as well as from a genuine class dependence, so we test the dependence statistically for each dataset and detector. We fit the per-seed effects with the linear model
\begin{equation}
\Delta_{o,c,k}=\mu+\alpha_k+\beta_o+\gamma_c+(\alpha\gamma)_{k,c}+(\beta\gamma)_{o,c}+e_{o,c,k}
\label{eq:anova}
\end{equation}
and compare it, with an $F$-test, against the reduced model without the edit$\times$class term $(\beta\gamma)_{o,c}$. The seed$\times$class term is required because all effects for a given seed and class are computed against the same base branch and therefore share its noise; the term absorbs this shared component exactly. A significant edit$\times$class term means that the effect of an edit depends on the class beyond what the average effect of the edit and the average sensitivity of the class explain. To locate the edits responsible, we then test each edit separately, comparing the model $\Delta_{o,c,k}=\mu+\alpha_k+\gamma_c+e_{c,k}$ with the same model without the class term, and correct the resulting $p$-values across edits with the Holm procedure. We conclude that classes respond differently to the same edits when the edit$\times$class term is significant at the 5\% level, or when an edit has an effect that is reliably positive for one class and reliably negative for another \todo{define ``reliably'': e.g.\ 95\% interval over seeds excluding zero}.

\subsection{Agreement Between Short and Full Training}
\label{sec:measure_proxy}

The short continuations measure the effect of an edit early in training, whereas CUA and the compared policies train for the full schedule. To check that the short runs rank edits as full training would, we train the base recipe and four edits for the full schedule on RDD2022 with Faster R-CNN and one seed, choosing the edits to span the ranking of the edit-wise study: the best, the worst and two from the middle. Agreement is measured by the Spearman rank correlation between the seed-averaged short-run effects $\bar\Delta(o,c)$ and the full-run effects over the shared edit--class cells, with a 95\% bootstrap interval from 2,000 resamples of the cells. If the correlation is below 0.5, we treat the short-run effects as indicative of direction only, and the measured-effect policy uses full-run effects wherever they are available.

\subsection{Class-Targeted Study}
\label{sec:measure_targeted}

The edit-wise study applies each edit to the whole image, so the effect it measures for a class combines two routes: the edit changes the appearance of that class's objects, and it changes the other objects and the background, which can influence the class through the parameters that all classes share. The class-targeted study separates the two routes by applying an edit to one class at a time. For a target class $c$, the edit is rendered only inside the class-$c$ boxes,
\begin{equation}
\tilde x=x+M_c\odot\bigl(T_o(x)-x\bigr),
\label{eq:targeted}
\end{equation}
where $M_c$ is the union mask of the class-$c$ boxes, feathered over 8 pixels at its border, and $\odot$ denotes pixelwise multiplication. The edit is applied with probability 0.5 to training images that contain $c$, within the same warm-start protocol as the edit-wise study \todo{confirm: same warm checkpoint and 3-epoch continuation}. The change in the AP50 of class $c$ is the direct effect of the edit, and the change for every other class $c'\neq c$ is its spill-over; collecting both over all target classes yields a direct/spill-over matrix for each edit. The study uses RDD2022 and dacl10k with Faster R-CNN and three seeds, and the three edits per dataset whose effects differ most across classes in the edit-wise study \todo{selection criterion, e.g.\ largest per-edit class $F$ statistic}. Geometric edits are excluded, because they cannot be confined to the boxes of one class without moving the boxes themselves. A local edit also creates a seam at the mask border that the detector could learn to exploit; to separate this artefact from the effect of the edit, one control run applies an edit to a random object-free region of the same area.

\section{Experimental Setup}
\label{sec:setup}

\subsection{Datasets}
\label{sec:setup_data}

We evaluate on three domains that differ in how their classes are defined and how often they co-occur (Table~\ref{tab:datasets}). RDD2022~\cite{arya2024rdd} contains road images from several countries with four damage classes: longitudinal cracks (D00), transverse cracks (D10), alligator cracks (D20) and potholes (D40). We use a fixed split of 26,863 training, 5,760 validation and 5,760 test images, in which exact duplicates and perceptually similar images are always assigned to the same partition. dacl10k~\cite{flotzinger2024dacl10k} contains close-up bridge-inspection images with polygon annotations of damage and bridge components. Because the official download is no longer available, we use its public FiftyOne export \todo{cite: Voxel51/dacl10k mirror on Hugging Face}, which contains 8,922 images under a CC BY-NC 4.0 licence. Its official test images carry no labels and are discarded, so that unlabelled images are never treated as damage-free. We keep the damage classes only, convert each annotated polygon into one box spanning all of its rings and clipped to the image, and drop degenerate polygons. Since the official splits are not available with labels, we group the labelled images by near-duplicate hashing and divide the groups 70/15/15 into training, validation and test partitions. PASCAL VOC~\cite{everingham2010voc} provides 20 everyday object classes; we train on the VOC2007 and VOC2012 trainval sets, from which 10\% of the images are held out for validation, and test on VOC2007 test. Objects marked difficult are excluded from the training targets and treated as ignore regions in evaluation.

\begin{table}[t]
\caption{Datasets, Splits and Co-Occurrence of Classes}
\label{tab:datasets}
\centering
\footnotesize
\setlength{\tabcolsep}{3pt}
\begin{tabular}{@{}lcccccc@{}}
\toprule
\textbf{Dataset} & \textbf{Classes} & \textbf{Train} & \textbf{Ref.} & \textbf{Val} & \textbf{Test} & \textbf{Multi-class} \\
\midrule
RDD2022~\cite{arya2024rdd} & 4 & \todo{n} & \todo{n} & 5,760 & 5,760 & \todo{\%} \\
dacl10k~\cite{flotzinger2024dacl10k} & \todo{n} & \todo{n} & \todo{n} & \todo{n} & \todo{n} & \todo{\%} \\
PASCAL VOC~\cite{everingham2010voc} & 20 & \todo{n} & \todo{n} & \todo{n} & 4,952 & \todo{\%} \\
\bottomrule
\end{tabular}
\vspace{2pt}

\parbox{\columnwidth}{\footnotesize Train excludes the reference set (Ref.), which only the class-utility estimator uses. Multi-class: share of object-bearing training images that contain two or more classes.}
\end{table}

\subsection{Data Gates and Reference Set}
\label{sec:setup_gates}

Three checks with fixed thresholds run once on each dataset before any training. First, because CUA addresses images that contain several classes, a dataset is used only if at least 20\% of its object-bearing training images contain two or more classes. Second, a dacl10k class is kept only if it has at least 300 training boxes and no more than half of its boxes cover more than half of the image; the first condition removes classes too rare to evaluate reliably, and the second removes region-like classes that a box describes poorly. This gate uses the training partition only, and it excludes \todo{list of excluded dacl10k classes}. Third, we screen for near-duplicate images with 64-bit difference hashes at a Hamming distance of at most 4: in dacl10k each group of near duplicates is kept within one partition, and for PASCAL VOC, which keeps its official partitions, duplicates across partitions are counted and reported.

From each training partition we then remove a reference set of about 1,000 images, chosen in whole near-duplicate groups so that each class has at least 100 objects where possible. The reference set is excluded from optimization in every compared policy, including the base recipe, so that all policies train on identical data. The remaining partitions have fixed roles: the training images are used for optimization, the reference set only for the class-utility estimates, the validation set for the measurements of Section~\ref{sec:measurement}, for tuning and for checkpoint selection, and the test set for the final comparison, which is evaluated once after all settings are fixed.

\subsection{Detectors and Training}
\label{sec:setup_training}

We use two detector families: Faster R-CNN with a ResNet-50 feature pyramid backbone~\cite{ren2017faster} and \todo{second detector family: RetinaNet~\cite{lin2017focal} as planned, or YOLOv8 (\todo{cite})}. Both start from COCO-pretrained torchvision weights, with all backbone stages trainable and batch-normalization statistics frozen. Images are letterboxed to $640\times640$ pixels with padding value 114. Training uses stochastic gradient descent with momentum 0.9, weight decay $5\times10^{-4}$, a batch size of four and a base learning rate of 0.0025, reached by a linear warm-up over 1,000 iterations from $0.001$ times that value; the learning rate is reduced tenfold after 60\% of the epochs and a hundredfold after 80\%, and training uses mixed precision. At inference, detections are kept above a score of 0.001 after non-maximum suppression at IoU 0.5, with at most 100 per image. Each run is validated after every epoch, and the checkpoint with the highest validation macro AP50 is the one evaluated on the test set. Seeds $\{0,1,2\}$ are shared across policies and fix the initialization of the new heads, the order of the training images and the augmentation draws.

The number of epochs is fixed for each dataset before policies are compared, from a base-recipe run at a candidate budget of 20 epochs for RDD2022 and PASCAL VOC and 30 for dacl10k. A budget is accepted if validation AP50 rises by less than 0.3 points over its last 20\% of epochs; otherwise it is extended by 50\% and checked again. The resulting budgets are \todo{RDD2022: 20 or 30} epochs for RDD2022, \todo{n} for PASCAL VOC and \todo{n} for dacl10k.

The settings of CUA are tuned once and then frozen. Using validation macro AP50 on RDD2022 with Faster R-CNN and one seed, we select the estimation interval $K\in\{100,200,400\}$, the parameter subset $\theta_s$ (box head alone, or box head and top feature-pyramid level) and the exploration weight $\varepsilon\in\{0.05,0.1,0.2\}$; the selected values, \todo{final $K$, $\theta_s$, $\varepsilon$}, are used unchanged for all datasets and detectors.

\subsection{Augmentation Menu}
\label{sec:setup_menu}

Table~\ref{tab:menu} lists the edits and their strengths. Each edit has one strength range, from which a value is drawn every time the edit is applied. Geometric edits transform the boxes together with the image, with a rotated box replaced by the axis-aligned hull of its rotated corners; a draw that pushes a box out of the canvas is redrawn up to ten times, after which the image is left unedited. The validity of each edit is fixed per dataset before training. Vertical flips are excluded for RDD2022, whose images are taken from a vehicle with the road always below the horizon, and for PASCAL VOC, whose upside-down scenes are not label-preserving in practice. Rotation is limited to $\pm10^\circ$ on every dataset, because larger rotations of road images would turn longitudinal cracks into transverse ones. Copy-paste augmentation is not included, because RDD2022 provides no instance masks.

\begin{table}[t]
\caption{Augmentation Menu and Validity per Dataset}
\label{tab:menu}
\centering
\footnotesize
\setlength{\tabcolsep}{3pt}
\begin{tabularx}{\columnwidth}{@{}llYccc@{}}
\toprule
\textbf{Family} & \textbf{Edit} & \textbf{Strength} & \textbf{RDD} & \textbf{VOC} & \textbf{dacl} \\
\midrule
Geometric & Translation & up to $\pm10\%$ of the canvas per axis & $\checkmark$ & $\checkmark$ & $\checkmark$ \\
 & Rotation & $\pm10^\circ$ & $\checkmark$ & $\checkmark$ & $\checkmark$ \\
 & Vertical flip & -- & -- & -- & $\checkmark$ \\
Colour & Brightness & factor in $[0.7,1.3]$ & $\checkmark$ & $\checkmark$ & $\checkmark$ \\
 & Contrast & factor in $[0.7,1.3]$ & $\checkmark$ & $\checkmark$ & $\checkmark$ \\
 & Saturation & factor in $[0.5,1.5]$ & $\checkmark$ & $\checkmark$ & $\checkmark$ \\
 & Hue & shift up to $\pm0.05$ of the hue circle & $\checkmark$ & $\checkmark$ & $\checkmark$ \\
Image quality & Gaussian blur & $\sigma\in[0.5,1.5]$ px & $\checkmark$ & $\checkmark$ & $\checkmark$ \\
 & Gaussian noise & $\sigma=0.03$ of the intensity range & $\checkmark$ & $\checkmark$ & $\checkmark$ \\
 & JPEG compression & quality in $[30,60]$ & $\checkmark$ & $\checkmark$ & $\checkmark$ \\
 & Sharpening & factor in $[1.5,2.5]$ & $\checkmark$ & $\checkmark$ & $\checkmark$ \\
Occlusion & Random erasing & 2--10\% of the canvas, noise-filled, at most 50\% of any box & $\checkmark$ & $\checkmark$ & $\checkmark$ \\
Composition & Mosaic & $2\times2$ tiles at half scale; boxes less than 30\% visible dropped & $\checkmark$ & $\checkmark$ & $\checkmark$ \\
\midrule
\multicolumn{3}{@{}l}{Valid edits} & 12 & 12 & 13 \\
\bottomrule
\end{tabularx}
\end{table}

\subsection{Compared Policies}
\label{sec:setup_policies}

We compare five policies, all applied on top of the base recipe and all trained on the same data, schedule and seeds. The \emph{base recipe} applies no further edit. \emph{Uniform augmentation} samples one edit uniformly from $\mathcal{O}\cup\{o_0\}$ for every image, in the manner of TrivialAugment~\cite{muller2021trivial}, and measures the value of the menu without any choice. The \emph{measured-effect policy} uses the policy of Section~\ref{sec:method_policy} with the class utilities replaced by the measured effects $\bar\Delta(o,c)$ of Section~\ref{sec:measurement}, and their standard errors by the standard errors over seeds \todo{confirm}; it is fixed throughout training and shows what class-specific choices achieve when the effects are measured offline rather than estimated online. \emph{Shared-utility augmentation} (Sec.~\ref{sec:method_shared}) and \emph{CUA} complete the comparison.

Ablations on the validation sets of RDD2022 and dacl10k, with Faster R-CNN and three seeds, remove or replace one component of CUA at a time: the harm block of Eq.~\eqref{eq:harm}; the identity candidate $o_0$; equal class weights, replaced by weights $w_c$ inversely proportional to the current validation AP of each class; and the class utility itself, replaced by a tolerance-gated controller that adapts edits according to whether the detector still detects the edited objects rather than whether training on them helps. \todo{Decide whether to add an ablation that combines tolerance and benefit.}

\subsection{Metrics and Statistics}
\label{sec:setup_metrics}

The primary metric is test macro AP50: average precision at an IoU threshold of 0.5, computed with the COCO evaluation protocol and averaged over the classes that have at least one test object. Secondary metrics are AP50:95, per-class AP50, AP for small, medium and large objects, and, for RDD2022, the fraction of object-free test images that receive at least one detection above confidence \todo{threshold}. We also report the wall-clock overhead of the estimator.

Every policy is trained with the same three seeds, so all comparisons are paired, and we report the mean and standard deviation over seeds. For the difference between two policies, we compute two 95\% intervals: a paired image-level bootstrap with 1,000 resamples of the test images, in which both policies are evaluated on the same resample within each seed before pooling, and a paired $t$ interval over seeds. We claim a difference only when both intervals exclude zero. The primary comparison is CUA against shared-utility augmentation on test macro AP50; if its bootstrap interval spans zero and is wider than 1.0 AP50, both policies are trained with two further seeds. To test whether CUA estimates what it is meant to estimate, we compute the Spearman correlation between the class utilities averaged over training and the measured effects $\bar\Delta(o,c)$, with a bootstrap interval over edit--class cells.

\subsection{Implementation}
\label{sec:setup_impl}

All methods are implemented in PyTorch 2.3 with torchvision 0.18 and trained on single NVIDIA A10 GPUs. Every run uses the same evaluation code, based on pycocotools. The complete study used \todo{measured GPU-hours} GPU-hours, of which the class-utility estimator accounted for \todo{share}.

